\section{Stream Blanker}

El stream blanker es el mecanismo que permite interrumpir la emisión hacia el público
sin interrumpir el flujo de trabajo interno del realizador. Se implementa en
\texttt{voctocore/lib/streamblanker.py} como un elemento GStreamer adicional que se
interpone entre la mezcla y el puerto de salida de stream (puerto 15000).

\subsection{Estados del sistema}

El sistema define tres estados de emisión:

\begin{itemize}
  \item \textbf{LIVE:} se emite la mezcla activa producida por voctocore. El realizador
        controla en todo momento lo que ve el espectador.
  \item \textbf{PAUSE:} se emite un bucle de vídeo de pausa (típicamente una imagen
        animada con música) con su canal de audio asociado.
  \item \textbf{NOSTREAM:} se emite un bucle de fuera de emisión, indicando al
        espectador que la transmisión ha finalizado o aún no ha comenzado.
\end{itemize}

La transición entre estados se realiza modificando los pesos del compositor en tiempo
real, evitando cortes bruscos en el stream de salida.

\subsection{Interfaz de control}

Desde la interfaz gráfica, el operador dispone de tres botones en la barra de
herramientas (\texttt{voctogui/lib/toolbar/streamblank.py}) etiquetados LIVE, PAUSE
y NOSTREAM. Al activar cualquiera de ellos se envía el comando
\texttt{set\_stream\_blank <estado>} al núcleo mediante el protocolo TCP.

% --- Figura sugerida ---
% \begin{figure}[H]
%   \centering
%   \includegraphics[width=0.6\textwidth]{figuras/fsm_stream_blanker.png}
%   \caption{Máquina de estados del Stream Blanker (LIVE / PAUSE / NOSTREAM).}
%   \label{fig:fsm_stream_blanker}
% \end{figure}

% --- Figura sugerida ---
% \begin{figure}[H]
%   \centering
%   \includegraphics[width=0.5\textwidth]{figuras/botones_stream_blanker.png}
%   \caption{Botones LIVE, PAUSE y NOSTREAM en la barra de herramientas de voctogui.}
%   \label{fig:botones_stream_blanker}
% \end{figure}
